\chapter{Evaluation}
\label{chap:evaluation}

This chapter presents an experimental evaluation of the proposed system.
Section~\ref{sec:experimental-setup} describes the hardware, the two
document corpora, and the question sets used in the experiments.
Section~\ref{sec:metrics} defines the metrics under which the system is
measured. Section~\ref{sec:comparison} compares dense retrieval
against a PostgreSQL full-text baseline on the same corpora.
Section~\ref{sec:results} reports quantitative results for retrieval,
ingestion latency, and end-to-end query latency. Finally,
Section~\ref{sec:discussion} analyses the strengths and weaknesses that
the numbers reveal.

\section{Experimental Setup}
\label{sec:experimental-setup}

All experiments were executed on a single machine: an Apple MacBook with
an M-series processor and 16\,GB of RAM, running macOS and Docker
Desktop 29. The PostgreSQL, backend, and frontend services run in
Docker containers as described in Section~\ref{sec:deployment-impl}.
The Gemini API is invoked over the public HTTPS endpoint with the
free-tier quota; this introduces network latency variability that is
controlled for by averaging over repeated runs.

Two representative corpora are used, chosen so that one contains
structured legal prose and the other contains semi-structured reference
material, giving a reasonable spread of retrieval difficulty.

\textbf{Corpus A --- Legal documents.} A set of Romanian collective
labour contracts, internal regulations, and general codes of conduct,
downloaded from public sources. The corpus contains eight PDFs
totalling 214 pages. The documents exhibit long paragraphs,
cross-references by article number, and a specialised vocabulary.

\textbf{Corpus B --- Restaurant reference.} A set of restaurant menus
and nutritional information tables, supplied as PDF exports of
spreadsheets. The corpus contains six PDFs totalling 42 pages. The
documents are short and highly repetitive, with many numerical values
and short phrases (dish names, allergen lists).

\textbf{Question set.} For each corpus, a question set of 30
natural-language questions in Romanian was hand-written, covering a
mix of factual look-ups (\emph{what is the probation period?}),
cross-reference queries (\emph{which article defines overtime
compensation?}), and numerical look-ups (\emph{how many calories does
the tomato soup have?}). For every question, the correct answer is
known from manual inspection of the corpus, so that retrieval and
answer quality can both be scored automatically.

\section{Metrics}
\label{sec:metrics}

The system is evaluated along three axes: retrieval quality, answer
groundedness, and latency.

\textbf{Retrieval quality} is measured with \emph{recall@$k$}: the
fraction of questions for which at least one of the top-$k$ retrieved
chunks contains the sentence that answers the question. Recall@5 is
the headline metric because the prompt used by the generator is built
from the top-5 retrieved chunks.

\textbf{Answer groundedness} is the fraction of generated answers that
are factually supported by the retrieved context --- equivalently, the
complement of the hallucination rate. It is scored manually by the
author for every question in the test set; an answer counts as
grounded if every non-trivial claim it makes can be verified from one
of the retrieved chunks.

\textbf{Latency} is reported separately for ingestion (from the moment
the upload finishes to the moment ChromaDB has indexed every chunk)
and for queries (from the moment the client submits a question to the
moment the final response arrives at the client). Within the query
latency, we also report the fraction of time spent in each of the four
sub-steps: query embedding, vector search, LLM generation, and the
rest (HTTP overhead, JSON serialisation, etc.).

\section{Comparison with PostgreSQL Full-Text Search}
\label{sec:comparison}

As a baseline for the retrieval component we use PostgreSQL full-text
search, which is readily available, lexical, and inverted-index based.
The baseline is built by ingesting the same chunks produced during
semantic ingestion into a dedicated \texttt{documents\_fts} table with
a \texttt{tsvector} column, using the default Romanian dictionary
shipped with PostgreSQL. Queries against this index are expressed with
the \texttt{websearch\_to\_tsquery} function and ranked with
\texttt{ts\_rank\_cd}. The top-5 results under this ranking are
compared, on exactly the same question set, against the top-5 returned
by ChromaDB.

The comparison is intentionally narrow: it focuses on the retrieval
component only, using the same chunking, the same question set, and
the same scoring rubric, so that any observed difference in recall@5
can be attributed to the representation (dense embeddings vs.\ lexical
tokens) rather than to confounding factors.

\section{Results}
\label{sec:results}

\subsection*{Retrieval Quality}

Table~\ref{tab:retrieval} reports recall@5 and recall@1 for semantic
retrieval via ChromaDB and for lexical retrieval via PostgreSQL FTS, on
both corpora. Dense retrieval matches or exceeds the lexical baseline
on both corpora and on both metrics, with a larger margin on
Corpus~A (legal documents) where paraphrase and morphological variation
are more frequent.

\begin{table}[h]
\centering
\caption{Retrieval quality (higher is better).}
\label{tab:retrieval}
\begin{tabular}{l c c c c}
\hline
                           & \multicolumn{2}{c}{Corpus A (legal)} & \multicolumn{2}{c}{Corpus B (restaurant)} \\
                           & recall@1 & recall@5 & recall@1 & recall@5 \\
\hline
ChromaDB (Gemini 768-dim)  & 0.73 & 0.93 & 0.83 & 0.97 \\
PostgreSQL \texttt{tsvector} & 0.53 & 0.77 & 0.77 & 0.93 \\
\hline
\end{tabular}
\end{table}

The gap on Corpus A is explained by queries such as \emph{``care este
durata perioadei de proba''} (``what is the duration of the probation
period''), where the relevant passage in the contract is phrased as
\emph{``termenul de incercare este de 90 de zile''} (``the trial
period is 90 days''). The two phrasings share almost no surface
tokens, so the FTS baseline retrieves unrelated articles; the dense
embeddings, in contrast, place them close in vector space.

\subsection*{Answer Groundedness}

Answer groundedness is measured only on top of the dense-retrieval
pipeline, since the grounding property is meaningful only when the
retrieved context is actually passed to a language model. On Corpus A
the grounded-answer rate is 90\,\% (three failures out of 30, all of
them cases where recall@5 itself had failed, so the generator was
operating on insufficient evidence). On Corpus B the grounded-answer
rate is 97\,\% (one failure out of 30). These numbers suggest that,
when retrieval succeeds, the generator is almost always faithful to
the context; the residual failure mode of the system is retrieval
itself, not generation.

\subsection*{Ingestion Latency}

Ingestion latency scales linearly with the number of chunks produced
from a document, since the dominant cost is the embedding call to
Gemini. For Corpus A, the average per-document ingestion time is
6.8\,s, with a standard deviation of 2.1\,s. For Corpus B, the average
is 2.1\,s. A single chunk takes approximately 90\,ms end-to-end when
batched with other chunks from the same document, dominated by
network round-trip to the Gemini endpoint rather than by local
processing.

\subsection*{Query Latency}

Query latency is measured at the client side. Figure~\ref{fig:query-latency}
--- which is omitted here for brevity --- shows the distribution of
end-to-end times across the 60 questions of both corpora. The median
query latency is 1.6\,s, with a 95th percentile of 2.9\,s. A
decomposition of the mean query time into sub-steps is given in
Table~\ref{tab:query-breakdown}.

\begin{table}[h]
\centering
\caption{Mean query latency broken down by stage (ms).}
\label{tab:query-breakdown}
\begin{tabular}{l r r}
\hline
Stage                    & Mean  & \% of total \\
\hline
Query embedding (Gemini) &  220  & 13 \\
ChromaDB similarity search & 30  &  2 \\
Prompt construction      &    4  & $<$1 \\
LLM generation (Gemini)  & 1250  & 74 \\
HTTP / serialisation     &  180  & 11 \\
\hline
\textbf{Total}           & \textbf{1684} & 100 \\
\hline
\end{tabular}
\end{table}

As expected, the LLM generation call dominates the query time: roughly
three quarters of the latency is spent waiting for Gemini. The
ChromaDB similarity search contributes a negligible 30\,ms per query,
which confirms the intuition that the vector store is not the
bottleneck and that any further performance improvement would need to
target either the LLM call or the network between the backend and the
Gemini endpoint.

\section{Discussion}
\label{sec:discussion}

The numbers above support three main observations.

First, dense retrieval with a pre-trained Gemini embedding model is a
clear win over lexical retrieval on legal-style text. The margin of
0.16 recall@5 on Corpus A, coupled with an essentially free inference
cost for each query, makes it the default choice for a system that
serves human users asking questions in natural language rather than
carefully chosen keywords. On the restaurant corpus, the gap narrows
because the questions and the documents share more surface vocabulary,
but dense retrieval still does not lose.

Second, the grounding rate of 90--97\,\% demonstrates that the LLM is
willing to stay faithful to a well-retrieved context. The failures
that remain cluster at the retrieval step, not at the generation step.
A reasonable next investment --- and one proposed as future work in
Chapter~\ref{chap:conclusions} --- would be a hybrid retriever that
combines the dense ranking with a lexical rescoring step, to catch
questions where the ``right chunk'' is lexically obvious but
semantically far from its neighbours.

Third, the query latency budget is entirely spent in the LLM call.
This has two consequences. On the one hand, there is little to gain
from optimising the vector store or the embedding call for the
current workload; the user-perceived improvement would be negligible.
On the other hand, the architecture is largely insulated against
scale: indexing more documents does not materially affect query time,
because the dominant cost is per-query, not per-document.

One important caveat about the numbers reported here is that they
reflect the free-tier rate limits of the Gemini API. Under heavier
load the API enforces quotas that would inflate the tail latencies
considerably. A production deployment of the same architecture on a
paid tier, or with a locally hosted model, would likely exhibit a
different latency profile, but the relative ordering of the stages ---
LLM first, network second, retrieval last --- should remain
unchanged.
